\chapter{Plan}\label{ch:plan}
In part two of the thesis, we will complete the implementation of a WasmGC backend and runtime system capable of running a set of standard Haskell benchmarks. We will then measure and evaluate the performance of the generated Wasm code, and develop a small demo application to showcase the capabilities of the WasmGC backend.

In the first ten weeks of part two, we will focus on implementing the WasmGC backend and runtime system. The following five weeks will be dedicated to setting up the benchmarks, measuring tooling and performance evaluation. The five weeks after that will be used to develop the demo application and/or implement optimizations, depending on the results of the performance evaluation and the time remaining (if the performance evaluation shows that the WasmGC backend is not performing well, we can prioritize optimizations to improve its performance before developing the demo application). The last four weeks will be used to write auxiliary sections of the thesis, prepare the defense, and catch up on any remaining work (in case of delays). If time permits, we can also explore additional features such as lightweight threads and JavaScript interoperability.

In case of major delays during step one and two, we can adjust the plan to not include the demo application and/or optimizations, and instead focus solely on the implementation of the WasmGC backend and runtime system and the performance evaluation. This would allow us to still have a complete implementation of the WasmGC backend and runtime system and a performance evaluation, which are the core contributions of this thesis.

\paragraph{Runtime Support for Necessary Features} We will not implement runtime support for all primitives MicroHs can generate, but instead we will define a new simplified prelude that is optimized for generating as few different types of combinators and primitives as possible. This prelude would provide the necessary functionality for the benchmarks while minimizing the number of combinators and primitives that need to be implemented in the runtime. We acknowledge that this approach may affect the performance of the generated code, but it allows us to focus on the core features of the WasmGC backend and runtime system without getting bogged down in implementing a large number of primitives. If time permits, we can explore adding support for more primitives in the runtime system to improve the performance of the generated code.

\paragraph{Performance Evaluation} With our performance evaluation we aim to answer the followig questions:
\begin{itemize}
    \item How does the performance of the WasmGC backend compare to MicroHs's existing C backend?
    \item How does the performance of the WasmGC backend compare to the MicroHs Wasm backend (through emscripten)?
    \item How does the performance of the WasmGC backend compare to the GHC Wasm backend?
    \item What performance improvements can be achieved through optimizations such as using a Wasm-to-Wasm optimizer like Binaryen?
\end{itemize}
Question one will give us insight into the performance impact of using WasmGC when compared to a traditional native backend. Based on this we can discuss if the performance of WasmGC is sufficient for real-world applications taking into account the benefits of using WasmGC such as portability. Question two will allow us to see if using WasmGC provides performance benefits compared to Wasm without GC, this will help compiler developers decide if it is worth it to target WasmGC instead of traditional Wasm. Question three will help us understand how well the WasmGC backend performs in relation to a mainstream Haskell compiler, this can give us insight into the potential of WasmGC for running real-world Haskell applications. To make this a fair comparison (as the runtime systems and optimizations performed by MicroHs and GHC differ significantly) we compare the relative performance of both backends to their respective native backends (i.e., MicroHs C backend and GHC's native backend). Question four will give us insight into how well the generated Wasm code can be optimized and how this compares to the optimizability of traditional Wasm.

To answer these questions we will run a selected set of benchmarks from the Nofib benchmark suite~\cite{nofib}, which is the defacto standard benchmark suite for Haskell compilers. NoFib consists of four sets of benchmark programs: imaginary, spectral, real, and shootout. We will not be able to run the real and shootout sets of benchmarks as they rely on features such as concurrency and foreign function interfaces that will not be implemented in the first version of our WasmGC backend. The imaginary and spectral sets of benchmarks consists of smaller algorithmic programs that do not rely on such features. We will start by running the imaginary benchmarks as they are the simplest benchmarks in terms of features used and then we will move on to the spectral benchmarks if time permits. 

Since we are comparing the performance of Wasm binaries to native binaries, we need to ensure that our measurements are fair and consistent across different environments (i.e. Haskell, C, and Wasm/JavaScript). To ensure this, we will perform research in the accuracy and consistency of different timing and measurement functions available in these environments and select those that are most similar in terms of accuracy and consistency.

\paragraph{Demo Application} We will develop a small demo application to showcase the capabilities of the WasmGC backend. This could for example be a simple version of the Haskell Playground~\cite{haskellplayground} that runs completely in the browser and allows users to write and execute Haskell code using the WasmGC backend. This would allow us to demonstrate the practical usability of the WasmGC backend and its potential for real-world applications.

\paragraph{Optimizations} There are multiple potential optimizations we can explore to improve the performance of the WasmGC backend. Some examples are as follows (others may be identified later on during the performance evaluation):
\begin{itemize}
    \item Removing recursion from the implementation of strict primitives, this can potentially improve performance as shown in the C runtime system of MicroHs. This runtime system makes use of tags which are put on the stack to indicate whether a value is being evaluated to be used by a strict primitive or not, checking for these tags on the stack makes it possible to not have to recurse. Results from the C runtime system show that this can lead to a $3\%$ performance increase or decrease depending on the hardware it is run on.
    \item Testing different stack allocation strategies for the LAS to optimize memory usage and performance. We cannot statically determine the necessary stack size for the LAS, so we currently allocate a fixed size array of $100000$ elements for the LAS. This may use more memory than necessary for some programs, while for other programs it may not be sufficient. We can explore different strategies for dynamically resizing the LAS or we can try using the Wasm value stack instead of a separate array for the LAS, this however will require very careful management of the value stack.
    \item Removing unused combinators and primitives from the runtime when a program is compiled. This has two potential benefits, first it can reduce the size of the generated Wasm binary which can lead to faster load times and second it leads to less branches (and thus less integer comparisons to check which combinator/primitive needs to be reduced) in the code responsible for reducing combinators and primitives which can lead to improved performance.
\end{itemize}

\paragraph{Lightweight Threads} To provide concurrency support in the runtime system we can implement lightweight threads using stack switching~\cite{stackswitching}. Every thread would be represented by its own stack and the program could switch between these stacks to achieve concurrency. This would involve implementing a scheduler to manage the threads and ensure that they are executed fairly.

\paragraph{JavaScript Interoperability} Since we are targeting the web, adding support for the Haskell-JavaScript foreign function interface can enable Haskell code to interact with JavaScript code running in the same environment. This would potentially allow features such as DOM manipulation, event handling, and access to browser APIs from Haskell code, which are necessary for building real-world web applications.